\chapter{Introduction}\label{sec:introduction}\markright{\thechapter~Introduction}

- Global status report \cite{UNEP2025}

- urgency for more sustainable construction of buildings \cite{SIA390_1_2025}

- Gebäude nach Bauperiode \cite{BFS2024_Bauperiode}

- why slabs: \>50\% of a normal concretes building mass - main lever for optimization \cite{DBV_Heft50_2024}

- Importance of early stage design \cite{broyles2024equations} - limited to structural implications

- Why functionality, economics matter too (flat slabs) \cite{Ammann2024}

- Tools and basics are not present that enable A comprehensive and fact-based comparison \cite{Kaufmann2024_WhitePaper}

- Practice oriented manner(systems that are already market ready)
    - novel concrete forms with structurally-optimized concrete floors \cite{Wangler2019, Menna2020}
    - Contractors and manufacturers have to invest significantly in construction technology to fabricate optimized systems \cite{Caulfield2022}

- Objective: Assess the impact of employing a certain floor system on the sustainability of a multi-story building
- Approach: Further develop the existing slab configurator and use it to systematically compare different floor systems that are optimised for sustainability but fulfill the required design criteria regarding structural safety, deflection, vibration, fire safety and building physics, whilst also taking other requirements into account. 


\section{Sustainability assessment during early-stage design}
During the conceptual design phase of a construction project, the potential to reduce the carbon embodied in a structure is greatest. However, due to the large number of possible design scenarios, conducting a comprehensive sustainability assessment at this stage can quickly exceed the available time and financial resources. Existing research on early-stage sustainability assessment is largely based on empirical studies that are not directly related to structural engineering calculations \cite{broyles2024equations}.
Thus, there is a need for tools that can systematically evaluate a wide range of combinations of design parameters while simultaneously satisfying all relevant design requirements, including ultimate limit state (ULS), serviceability limit state (SLS), and fire resistance criteria. Such tools would enable designers to assess and compare the sustainability performance of structurally feasible solutions already during the early stages of the design process.